% =============================================================================
%  II. BACKGROUND AND RELATED WORK
% =============================================================================
\begin{frame}{II. Background and Related Work}
{\small Wang et al.~\cite{wang2025srpm} reviewed 91 real-time sepsis
prediction models and found:}

\vspace{1pt}
\begin{center}\tablestyle\footnotesize
\begin{tabular}{lcc}
\toprule
\textbf{Metric} & \textbf{Internal} & \textbf{External} \\
\midrule
Median AUROC         & 0.811              & 0.783 \\
Median Utility Score & \textbf{+0.381}    & \textbf{--0.164} \\
\bottomrule
\end{tabular}
\end{center}

\vspace{1pt}
{\small\textbf{Four key limitations of the conventional approach:}}
\vspace{-2pt}
\begin{enumerate}\itemsep0pt\small
  \item \textbf{The sepsis label can change.} Different Sepsis-3 criteria
        create very different patient groups
        (867--2,178 cases)~\cite{cohen2024subtle}.
  \item \textbf{The data split can give unrealistic results.} Performance
        drops significantly with temporal splits~\cite{guo2022temporal}.
  \item \textbf{AUROC does not show the full picture.} A model can maintain
        good AUROC while clinical usefulness
        decreases~\cite{wang2025srpm}.
  \item \textbf{The label may represent doctor decisions}, not patient
        deterioration~\cite{kamran2024evaluation}.
\end{enumerate}
\end{frame}
